\section{Tecnologie utilizzate}

\subsection{Backend}
\begin{itemize}
    \item \textbf{Python 3.10+} come linguaggio.
    \item \textbf{Flask 3.1} come framework HTTP per la definizione delle
    rotte, gestione di middleware e error handler.
    \item \textbf{Flask-SQLAlchemy 3.1} come ORM, con \emph{session pattern}
    legato al contesto della richiesta.
    \item \textbf{SQLite} come database embedded (file \code{lyrics.db}),
    scelto per portabilità ed assenza di configurazione esterna.
    \item \textbf{Flask-JWT-Extended 4.7} per emissione e verifica dei JWT
    (algoritmo HS256), con \emph{access token} (30 min) e \emph{refresh token}
    (7 giorni).
    \item \textbf{Flask-CORS 6.0} per la gestione delle policy
    \emph{Cross-Origin Resource Sharing}.
    \item \textbf{werkzeug.security} per l'hashing delle password con
    PBKDF2-SHA256 e salt randomico.
\end{itemize}

\subsection{Mobile/Web}
\begin{itemize}
    \item \textbf{React Native 0.81} con \textbf{Expo SDK 54} per
    cross-platform iOS, Android e Web.
    \item \textbf{expo-router 6} per il routing file-based.
    \item \textbf{TypeScript} per type safety end-to-end.
    \item \textbf{expo-secure-store} (opzionale) per persistere il JWT in
    Keychain (iOS) o Keystore (Android). Su web il fallback è
    \code{localStorage}.
\end{itemize}

\subsection{Protocolli di comunicazione}
\begin{itemize}
    \item \textbf{HTTP/HTTPS} con payload \textbf{JSON} (RFC 8259).
    \item \textbf{REST} come stile architetturale.
    \item \textbf{JWT} (RFC 7519) per l'autenticazione stateless.
    \item \textbf{Proof of Work} ispirato a \emph{hashcash}: SHA-256 con
    target adattivo.
    \item Formato dei testi sincronizzati: \textbf{LRC} standard
    (\code{[mm:ss.xx] testo}).
\end{itemize}

\subsection{Strumenti di sviluppo e test}
\begin{itemize}
    \item \textbf{Git} per il version control.
    \item \textbf{Overleaf} / \LaTeX{} per la documentazione.
    \item \textbf{Postman} / \code{curl} per l'esplorazione manuale dell'API.
    \item Script di smoke test in Python che esercitano gli endpoint via
    \code{app.test\_client()} di Flask.
\end{itemize}
